% Bibitems used by the imprimitive and semiregular CM sections.
% Include inside the main paper's thebibliography environment.
% Native URLs and exact consulted locators are also in SOURCES.md.

\bibitem{IM-LehrerTaylor}
G.~I. Lehrer and D.~E. Taylor,
\emph{Unitary Reflection Groups}, Cambridge University Press, 2009,
Chapter~2, \S\S1--3, pp.~23--27, especially Proposition~2.10;
authors' errata, 5 February 2025,
\url{https://www.maths.usyd.edu.au/u/don/papers/errata-uggr.pdf}.

\bibitem{IM-LewisWang}
J.~B. Lewis and J.~Wang,
The Hurwitz action in complex reflection groups,
\emph{Combinatorial Theory} \textbf{2}(1) (2022), \S2.1,
\url{https://escholarship.org/uc/item/4xg2m191}.

\bibitem{IM-LangeOrtega}
H.~Lange and A.~Ortega,
Prym varieties of cyclic coverings,
\emph{Geometriae Dedicata} \textbf{150} (2011), 391--403,
doi:10.1007/s10711-010-9512-9.
Proposition~4.1 and proof, author version pp.~7--10,
\url{https://arxiv.org/abs/0805.1020}.

\bibitem{IM-FarbMargalit}
B.~Farb and D.~Margalit,
\emph{A Primer on Mapping Class Groups}, Princeton University Press, 2012.
Proposition~6.2, p.~166; Theorem~6.4, p.~170;
Theorem~4.6 and Fact~4.7, pp.~97--98.
Authors' publication page:
\url{https://www.math.uchicago.edu/~farb/books.html}.

\bibitem{IM-StacksDeform}
The Stacks Project authors,
\emph{The Stacks Project}, Deformation Theory, \S93.10,
Lemma~93.10.3, tag~0E3V,
\url{https://stacks.math.columbia.edu/tag/0E3V}.
Used as background only; the morphism-deformation argument is proved in
D-IM-07.

\bibitem{SR-BuchweitzFlenner2003}
R.-O. Buchweitz and H.~Flenner,
A semiregularity map for modules and applications to deformations,
\emph{Compositio Mathematica} \textbf{137} (2003), 135--210,
doi:10.1023/A:1023999012081.
\S4.1 and Proposition~4.2, pp.~166--167.

\bibitem{SR-BuchweitzFlenner2008}
R.-O. Buchweitz and H.~Flenner,
The global decomposition theorem for Hochschild (co-)homology of singular
spaces via the Atiyah--Chern character,
\emph{Advances in Mathematics} \textbf{217} (2008), 243--281,
doi:10.1016/j.aim.2007.06.013.
Proposition~6.2.1 and Theorem~6.4.2 in the author version,
pp.~29 and~31,
\url{https://arxiv.org/abs/math/0606730}.

\bibitem{SR-Mukai1981}
S.~Mukai,
Duality between $D(X)$ and $D(\widehat X)$ with its application to
Picard sheaves,
\emph{Nagoya Mathematical Journal} \textbf{81} (1981), 153--175,
doi:10.1017/S002776300001922X.
Theorem~2.2, p.~156, and equation~(3.1), p.~158.

\bibitem{SR-Beauville2014}
A.~Beauville,
Some surfaces with maximal Picard number,
\emph{Journal de l'\'Ecole polytechnique---Math\'ematiques}
\textbf{1} (2014), 101--116,
doi:10.5802/jep.5.
\S4, Lemma~2, p.~106, and Example~3, p.~107,
\url{https://www.numdam.org/article/JEP_2014__1__101_0.pdf}.

\bibitem{SR-Perry2026}
A.~Perry,
\emph{The semiregularity theorem for equivariant noncommutative varieties},
arXiv:2604.00511v2, 24 June 2026 (preprint).
Theorem~1.2, p.~3; Remarks~2.3 and~2.5, p.~5;
Definition~2.6 and Lemma~2.7, pp.~5--6,
\url{https://arxiv.org/abs/2604.00511}.

\bibitem{SR-BirkenhakeLange}
C.~Birkenhake and H.~Lange,
\emph{Complex Abelian Varieties}, second edition, Springer, 2004,
Proposition~11.2.1 (Poincar\'e formula),
\url{https://link.springer.com/book/10.1007/978-3-662-06307-1}.
The proposition is a classical input; the original book text was not
directly retrieved in this run. An independently consulted primary use
is Schoen~\cite{SR-Schoen1988}, p.~25.

\bibitem{SR-BauerSzemberg}
T.~Bauer and T.~Szemberg,
On tensor products of ample line bundles on abelian varieties,
\emph{Mathematische Zeitschrift} \textbf{223} (1996), 79--85,
Theorem~1.1(b).
The author paper's theorem and proof were consulted; three equal ample
factors give the very ampleness of $3\Theta$ used in D-SR-01.

\bibitem{SR-Schoen1988}
C.~Schoen,
Hodge classes on self-products of a variety with an automorphism,
\emph{Compositio Mathematica} \textbf{65} (1988), 3--32.
The simple-branch definition and Theorem~2.0, p.~11;
Corollary~3.1 and proof, pp.~24--25,
\url{https://www.numdam.org/item/CM_1988__65_1_3_0.pdf}.

\bibitem{IM:Looijenga1997}
E.~Looijenga,
Prym representations of mapping class groups,
\emph{Geometriae Dedicata} \textbf{64} (1997), 69--83,
doi:10.1023/A:1004909416648.
Proposition~2.2 and Theorem~2.4, author version p.~5;
the hyperbolic-basis proof, pp.~8--9, and monodromy proof for
genus at least three, pp.~10--12.
Author's PostScript: \url{https://webspace.science.uu.nl/~looij101/pr2.ps}.
